\paragraph{对角闭包与单次对角观测}
在定理 \ref{thm:pom-maxent-markov-diagonal-plus-rankone-spectrum} 的“对角 \(+\) 秩一”规范下，
最优核的非对角部分具有严格秩一结构；并且其对角保持率向量在代数意义下已足以回收对偶参数、失真水平、平稳分布与全核。

\begin{proposition}[对角闭包：非对角核的秩一乘积形]\label{prop:pom-diagonal-rate-diagonal-closure-rankone}
\leanverified{paper\_pom\_diagonal\_rate\_diagonal\_closure\_rankone}
取全支撑分布 \(w\in\Delta(X)\) 并令 \(\delta\in(0,\delta_0)\)，其中 \(\delta_0:=1-\sum_{x\in X}w(x)^2\)。
沿用定理 \ref{thm:pom-maxent-markov-diagonal-plus-rankone-spectrum} 与定理 \ref{thm:pom-maxent-markov-laguerre-secular-spectrum} 的记号，
令 \(P^\star_\delta,K^\star_\delta,u_\delta,\kappa_\delta,A_\delta\) 与
\(
t_x=\kappa_\delta+A_\delta/u_\delta(x)
\)
如上定义，并记对角保持率
\[
p_\delta(x):=K^\star_\delta(x\mid x)\in(0,1)\qquad(x\in X).
\]
则对任意 \(x,y\in X\) 有显式核公式
\[
K^\star_\delta(y\mid x)=
\begin{cases}
\dfrac{1+\kappa_\delta}{t_x},&y=x,\\[8pt]
\dfrac{t_x-\kappa_\delta}{t_x}\cdot\dfrac{1}{t_y-\kappa_\delta},&y\ne x,
\end{cases}
\qquad
\sum_{y\in X}\frac{1}{t_y-\kappa_\delta}=1.
\]
从而非对角部分严格为秩一乘积：对 \(y\ne x\) 有 \(K^\star_\delta(y\mid x)=a_\delta(x)b_\delta(y)\)，其中
\[
a_\delta(x):=\frac{t_x-\kappa_\delta}{t_x},
\qquad
b_\delta(y):=\frac{1}{t_y-\kappa_\delta}.
\]
\end{proposition}
\begin{proof}[证明要点]
由结论 \ref{cor:pom-maxent-markov-optimal-coupling-latent-binary-mixture} 的定义，
\(
w(x)=u_\delta(x)\bigl(A_\delta+\kappa_\delta u_\delta(x)\bigr)
=t_xu_\delta(x)^2
\)，
故
\(
p_\delta(x)=P^\star_\delta(x,x)/w(x)=(1+\kappa_\delta)u_\delta(x)^2/w(x)=(1+\kappa_\delta)/t_x
\)。
又由 \(t_x-\kappa_\delta=A_\delta/u_\delta(x)\) 得
\(
b_\delta(x)=1/(t_x-\kappa_\delta)=u_\delta(x)/A_\delta
\)，
从而 \(\sum_y b_\delta(y)=1\)。
最后，对 \(y\ne x\)，
\[
K^\star_\delta(y\mid x)=\frac{P^\star_\delta(x,y)}{w(x)}=\frac{u_\delta(x)u_\delta(y)}{t_xu_\delta(x)^2}
=\frac{u_\delta(y)}{u_\delta(x)t_x}
=\frac{A_\delta}{u_\delta(x)t_x}\cdot \frac{u_\delta(y)}{A_\delta}
=\frac{t_x-\kappa_\delta}{t_x}\cdot\frac{1}{t_y-\kappa_\delta}.
\]
证毕。
\end{proof}

\paragraph{Sturm--Cauchy 本征基、全 resolvent 闭式与刷新接受强平稳时刻}
在最大熵最优核 \(K^\star_\delta\) 的“对角 \(+\) 秩一”结构下，本征向量可由 Cauchy 型有理函数显式写出；
同时 \((I-sK^\star_\delta)^{-1}\) 的矩元由单一 Laguerre 世俗多项式 \(\mathcal Q_\delta\) 控制。
该结构还允许仅依赖刷新样本构造强平稳停止时间，并给出其有理概率母函数与可计算尾界。
\par

\begin{theorem}[Sturm--Cauchy 本征基：\texorpdfstring{$(t_x-\kappa_\delta)/(t_x-z)$}{(t-κ)/(t-z)} 给出显式本征函数]\label{thm:pom-diagonal-rate-sturm-cauchy-eigenbasis}
\leanverified{paper\_pom\_diagonal\_rate\_sturm\_cauchy\_eigenbasis}
取全支撑分布 \(w\in\Delta(X)\)（\(n:=|X|\ge 2\)）并固定 \(\delta\in(0,\delta_0)\)。
沿用定理 \ref{thm:pom-maxent-markov-laguerre-secular-spectrum} 与命题 \ref{prop:pom-diagonal-rate-diagonal-closure-rankone} 的记号：
\(K^\star_\delta\) 为唯一最优核，\(\kappa_\delta>0\)，\(\{t_x\}_{x\in X}\subset(\kappa_\delta,\infty)\)，并记
\[
\mathcal P_\delta(z):=\prod_{x\in X}(t_x-z),
\qquad
\mathcal Q_\delta(z):=z\,\mathcal P_\delta'(z)+\kappa_\delta\,\mathcal P_\delta(z).
\]
假设 \(t_x\) 两两不同，并令 \(\{z_i\}_{i=1}^{n}\) 为 \(\mathcal Q_\delta\) 的全部零点（按升序排列），并定义
\[
\lambda_i:=\frac{\kappa_\delta}{z_i},
\qquad
\phi_i(x):=\frac{t_x-\kappa_\delta}{t_x-z_i}\qquad(x\in X).
\]
则对每个 \(i=1,\dots,n\) 有 \(\phi_i\not\equiv 0\) 且满足本征关系
\[
(K^\star_\delta\phi_i)(x)=\lambda_i\,\phi_i(x)\qquad(\forall x\in X).
\]
其中 \(z_1=\kappa_\delta\)、\(\lambda_1=1\) 且 \(\phi_1\equiv 1\)。
并且当 \(i\ne j\) 时成立严格正交性
\[
\sum_{x\in X}w(x)\,\phi_i(x)\phi_j(x)=0.
\]
\end{theorem}
\begin{proof}
由命题 \ref{prop:pom-diagonal-rate-diagonal-closure-rankone}，最优核可写为秩一刷新分解
\[
K^\star_\delta(y\mid x)=(1-r_\delta(x))\mathbf 1_{\{y=x\}}+r_\delta(x)\pi_\delta(y),
\qquad
r_\delta(x)=\frac{t_x-\kappa_\delta}{t_x},
\quad
\pi_\delta(y)=\frac{1}{t_y-\kappa_\delta}.
\]
因此对任意函数 \(\phi\) 有
\[
(K^\star_\delta\phi)(x)=\frac{\kappa_\delta}{t_x}\phi(x)+\frac{t_x-\kappa_\delta}{t_x}\sum_{y\in X}\frac{\phi(y)}{t_y-\kappa_\delta}.
\]
取 \(\phi=\phi_i\)，并注意 \(\phi_i(y)/(t_y-\kappa_\delta)=1/(t_y-z_i)\)。
由定理 \ref{thm:pom-maxent-markov-laguerre-secular-spectrum} 的世俗方程（等价于 \(\mathcal Q_\delta(z_i)=0\)）
\[
\sum_{y\in X}\frac{1}{t_y-z_i}=\frac{\kappa_\delta}{z_i},
\]
得到
\[
(K^\star_\delta\phi_i)(x)
=\frac{\kappa_\delta}{t_x}\frac{t_x-\kappa_\delta}{t_x-z_i}
\frac{t_x-\kappa_\delta}{t_x}\cdot\frac{\kappa_\delta}{z_i}
=\frac{\kappa_\delta}{z_i}\cdot\frac{t_x-\kappa_\delta}{t_x-z_i}
=\lambda_i\phi_i(x),
\]
即本征关系。
由于 \(z_1=\kappa_\delta\)（同一定理），有 \(\phi_1(x)=(t_x-\kappa_\delta)/(t_x-\kappa_\delta)=1\)。
\par
再由定理 \ref{thm:pom-diagonal-rate-diagonal-statistics-complete} 的闭式
\(
w(x)=A_\delta^2\,t_x/(t_x-\kappa_\delta)^2
\)
（其中 \(A_\delta=\sum_{x}u_\delta(x)\)），可得
\[
w(x)\phi_i(x)\phi_j(x)=A_\delta^2\frac{t_x}{(t_x-z_i)(t_x-z_j)}.
\]
部分分式分解给出
\[
\frac{t_x}{(t_x-z_i)(t_x-z_j)}
=\frac{z_i}{z_i-z_j}\frac{1}{t_x-z_i}-\frac{z_j}{z_i-z_j}\frac{1}{t_x-z_j}.
\]
对 \(x\) 求和并用世俗方程 \(\sum_x (t_x-z_i)^{-1}=\kappa_\delta/z_i\)、\(\sum_x (t_x-z_j)^{-1}=\kappa_\delta/z_j\) 得总和为 \(0\)，从而正交性成立。证毕。
\end{proof}

\begin{proposition}[全 resolvent 矩元闭式：仅由 \texorpdfstring{$\mathcal Q_\delta(\kappa_\delta s)$}{Q(κs)} 控制]\label{prop:pom-diagonal-rate-resolvent-entry-closed-form}
\leanverified{paper\_pom\_diagonal\_rate\_resolvent\_entry\_closed\_form}
沿用定理 \ref{thm:pom-diagonal-rate-sturm-cauchy-eigenbasis} 的设定与记号。
对 \(|s|\) 充分小定义 resolvent 矩阵
\[
\mathcal G(s):=(I-sK^\star_\delta)^{-1},
\qquad
\mathcal G_{xy}(s):=\mathcal G(s)_{x,y}.
\]
则对任意 \(x,y\in X\) 有严格恒等式
\[
\boxed{\;
\mathcal G_{xy}(s)
=\frac{t_x}{t_x-\kappa_\delta s}\,\mathbf 1_{\{x=y\}}
\;+\;
\frac{s(t_x-\kappa_\delta)}{t_x-\kappa_\delta s}\cdot
\frac{t_y}{(t_y-\kappa_\delta)(t_y-\kappa_\delta s)}\cdot
\frac{\kappa_\delta\,\mathcal P_\delta(\kappa_\delta s)}{\mathcal Q_\delta(\kappa_\delta s)}\; }.
\]
\end{proposition}
\begin{proof}
由秩一刷新分解 \(K^\star_\delta=\mathrm{diag}(1-r_\delta)+r_\delta\,\pi_\delta^{\mathsf T}\)，令
\[
A(s):=I-s\,\mathrm{diag}(1-r_\delta)=\mathrm{diag}\!\Bigl(1-\frac{\kappa_\delta s}{t_x}\Bigr)_{x\in X}.
\]
则 \(I-sK^\star_\delta=A(s)-s\,r_\delta\,\pi_\delta^{\mathsf T}\)。
对秩一扰动应用 Sherman--Morrison 公式得到
\[
(I-sK^\star_\delta)^{-1}=A(s)^{-1}
\;+\;\frac{A(s)^{-1}(sr_\delta)\pi_\delta^{\mathsf T}A(s)^{-1}}{1-\pi_\delta^{\mathsf T}A(s)^{-1}(sr_\delta)}.
\]
由于 \(A(s)\) 为对角矩阵，有
\(
A(s)^{-1}_{xx}=t_x/(t_x-\kappa_\delta s)
\)。
并且 \(\pi_\delta(x)r_\delta(x)=1/t_x\)，从而
\[
\pi_\delta^{\mathsf T}A(s)^{-1}(sr_\delta)
=s\sum_{u\in X}\frac{1}{t_u-\kappa_\delta s}.
\]
令 \(z=\kappa_\delta s\)，并用对数导数恒等式
\(
\sum_{u}(t_u-z)^{-1}=-\mathcal P_\delta'(z)/\mathcal P_\delta(z)
\)
得
\[
1-\pi_\delta^{\mathsf T}A(s)^{-1}(sr_\delta)
=1-s\sum_{u}\frac{1}{t_u-\kappa_\delta s}
=\frac{\mathcal Q_\delta(\kappa_\delta s)}{\kappa_\delta\,\mathcal P_\delta(\kappa_\delta s)}.
\]
将以上表达代入并取 \((x,y)\) 元即可得到所示闭式。证毕。
\end{proof}

\begin{theorem}[谱投影、残数公式与精确转移闭式]\label{thm:pom-diagonal-rate-spectral-projector-residue-and-transition}
\leanverified{paper\_pom\_diagonal\_rate\_spectral\_projector\_residue\_and\_transition}
沿用定理 \ref{thm:pom-diagonal-rate-sturm-cauchy-eigenbasis} 的设定，并令
\[
\mathsf N_i:=\sum_{u\in X}w(u)\,\phi_i(u)^2
=\frac{\sum_{u\in X}\dfrac{t_u}{(t_u-z_i)^2}}{\sum_{u\in X}\dfrac{t_u}{(t_u-\kappa_\delta)^2}}
\qquad(i=1,\dots,n).
\]
定义核
\[
\Pi_i(x,y):=w(y)\,\frac{\phi_i(x)\phi_i(y)}{\mathsf N_i}\qquad(1\le i\le n).
\]
则 \(\Pi_i\) 为 \(K^\star_\delta\) 在 \(\ell^2(w)\) 上对应于特征值 \(\lambda_i\) 的谱投影，满足
\[
\Pi_i^2=\Pi_i,\qquad \Pi_i\Pi_j=0\ (i\ne j),\qquad \sum_{i=1}^{n}\Pi_i=I,
\]
并且对一切整数 \(m\ge 0\) 有精确展开
\[
\boxed{\;
(K^\star_\delta)^m(y\mid x)=\sum_{i=1}^{n}\lambda_i^m\,\Pi_i(x,y)\; }.
\]
此外，对 \(i\ge 2\) 还有等价的残数表达（取 \(s_i:=1/\lambda_i=z_i/\kappa_\delta\)）：
\[
\boxed{\;
\Pi_i(x,y)=-\lambda_i\,\mathrm{Res}_{s=s_i}\ \mathcal G_{xy}(s)\; } ,
\]
并且由命题 \ref{prop:pom-diagonal-rate-resolvent-entry-closed-form} 的闭式可化为
\[
\boxed{\;
\Pi_i(x,y)
=-\frac{\phi_i(x)}{t_y-\kappa_\delta}\cdot\frac{t_y}{t_y-z_i}\cdot
\frac{\mathcal P_\delta(z_i)}{\mathcal Q_\delta'(z_i)}\; } .
\]
\end{theorem}
\begin{proof}
由于 \(P^\star_\delta(x,y)=w(x)K^\star_\delta(y\mid x)\) 对称，\(K^\star_\delta\) 在 \(\ell^2(w)\) 上自伴随；
由定理 \ref{thm:pom-diagonal-rate-sturm-cauchy-eigenbasis}，\(\{\phi_i\}\) 给出 \(n\) 个两两正交的本征函数。
归一化 \(\psi_i:=\phi_i/\sqrt{\mathsf N_i}\) 得到 \(\ell^2(w)\) 的正交规范基，从而标准谱分解给出
\(
K^\star_\delta=\sum_i \lambda_i\,\Pi_i
\)
与
\(
(K^\star_\delta)^m=\sum_i \lambda_i^m\,\Pi_i
\)，即得到转移闭式。
\par
对残数公式，有限维谱分解给出
\(
(I-sK^\star_\delta)^{-1}=\sum_i \Pi_i/(1-s\lambda_i)
\)，故在 \(s=s_i=1/\lambda_i\) 的留数为 \(-\Pi_i/\lambda_i\)。
\par
最后，由命题 \ref{prop:pom-diagonal-rate-resolvent-entry-closed-form} 可知 \(\mathcal G_{xy}\) 的该极点仅来自因子
\(
\kappa_\delta\mathcal P_\delta(\kappa_\delta s)/\mathcal Q_\delta(\kappa_\delta s)
\)。
对 \(s_i=z_i/\kappa_\delta\) 处取留数得
\(
\mathrm{Res}_{s=s_i}\bigl(\kappa_\delta\mathcal P_\delta(\kappa_\delta s)/\mathcal Q_\delta(\kappa_\delta s)\bigr)
=\mathcal P_\delta(z_i)/\mathcal Q_\delta'(z_i)
\)，
并代入其余解析因子即可得到显式表达式。证毕。
\end{proof}

\paragraph{刷新接受强平稳停止时间与其有理母函数}
秩一刷新表示使刷新样本在每次刷新时刻均服从同一背景分布 \(\pi_\delta\)，从而可在刷新序列之上构造拒绝采样型强平稳停止时间。

\begin{theorem}[刷新接受强平稳停止时间]\label{thm:pom-diagonal-rate-refresh-accept-strong-stationary-time}
\leanverified{paper\_pom\_diagonal\_rate\_refresh\_accept\_strong\_stationary\_time}
沿用命题 \ref{prop:pom-diagonal-rate-diagonal-closure-rankone} 的秩一刷新表示
\(
K^\star_\delta=\mathrm{diag}(1-r_\delta)+r_\delta\,\pi_\delta^{\mathsf T}
\)。
令
\[
c:=\min_{x\in X}r_\delta(x)\in(0,1),
\qquad
\alpha(y):=\frac{c}{r_\delta(y)}\in(0,1]\qquad(y\in X).
\]
考虑如下耦合：每次触发刷新时独立采样 \(Y\sim\pi_\delta\)，并在给定 \(Y\) 后以条件概率 \(\alpha(Y)\) 决定是否停止；
若停止则以该次刷新落点作为停止状态。
记由起点 \(X_0=x\) 导出的首次停止时间为 \(T\)。
则 \(T\) 为 \(K^\star_\delta\) 的强平稳停止时间：对任意 \(m\ge 0\) 与任意 \(y\in X\)，有
\[
\PP_x(T=m,\ X_T=y)=\PP_x(T=m)\,w(y),
\]
从而 \(X_T\sim w\) 且 \(X_T\) 与 \(T\) 独立。
\end{theorem}
\begin{proof}
在每次刷新时刻，候选落点 \(Y\) 的分布为 \(\pi_\delta\)，且接受事件在给定 \(Y\) 后以概率 \(\alpha(Y)\) 发生。
因此一次刷新轮内“被接受且落在 \(y\)”的概率与 \(\pi_\delta(y)\alpha(y)\) 成正比。
由定理 \ref{thm:pom-diagonal-rate-diagonal-statistics-complete}，
\(
w(y)\propto \pi_\delta(y)/r_\delta(y)
\)，
而 \(\pi_\delta(y)\alpha(y)=c\,\pi_\delta(y)/r_\delta(y)\)，故被接受的刷新落点严格服从 \(w\)。
\par
另一方面，刷新候选序列为独立同分布，接受—拒绝机制是标准拒绝采样：被接受样本的分布与“成功轮数”独立。
而停止时间 \(T\) 仅由成功轮数与此前各轮的停留长度决定（停留长度在给定轮落点后独立），故 \(X_T\) 与 \(T\) 独立。
将独立性与边缘分布合并即得联合因子化式。证毕。
\end{proof}

\begin{proposition}[仅依赖刷新样本的策略类：实现 \(w\) 的充要形式与最优性]\label{prop:pom-diagonal-rate-refresh-accept-uniqueness-optimality}
\leanverified{paper\_pom\_diagonal\_rate\_refresh\_accept\_uniqueness\_optimality}
在定理 \ref{thm:pom-diagonal-rate-refresh-accept-strong-stationary-time} 的“每次刷新给出候选 \(Y\sim\pi_\delta\)，并以某函数 \(a(Y)\in[0,1]\) 决定是否停止”的策略类中：
\begin{enumerate}
  \item 若要求停止状态分布恰为 \(w\)，则策略必且仅必形如
  \[
  a(y)=\frac{\theta}{r_\delta(y)}\qquad(\forall y\in X),
  \]
  其中常数 \(\theta\in(0,c]\) 且 \(c=\min_x r_\delta(x)\)。
  \item 在所有满足 \(X_T\sim w\) 的策略中，取 \(\theta=c\)（即 \(a(y)=c/r_\delta(y)\)）唯一同时最小化期望刷新轮数与期望停止时间。
\end{enumerate}
\end{proposition}
\begin{proof}
对任意策略 \(a\)，令 \(p:=\sum_y \pi_\delta(y)a(y)\) 为单轮成功概率。
则被接受的刷新落点分布为 \(\pi_\delta(y)a(y)/p\)。
要求其等于 \(w(y)\) 意味着存在常数 \(p\) 使 \(\pi_\delta(y)a(y)=p\,w(y)\) 对所有 \(y\) 成立。
由 \(w(y)\propto \pi_\delta(y)/r_\delta(y)\) 得 \(a(y)=\theta/r_\delta(y)\)，并由 \(a(y)\le 1\) 推出 \(\theta\le \min_y r_\delta(y)=c\)。
反向代入直接验证 \(X_T\sim w\)。
\par
在形式 \(a(y)=\theta/r_\delta(y)\) 下，成功概率为
\[
p(\theta)=\sum_y \pi_\delta(y)\frac{\theta}{r_\delta(y)}=\theta\sum_y\frac{\pi_\delta(y)}{r_\delta(y)}=\theta\,A_\delta^{-2},
\]
其中 \(A_\delta^{-2}=\sum_y t_y/(t_y-\kappa_\delta)^2\)（定理 \ref{thm:pom-diagonal-rate-diagonal-statistics-complete}）。
故刷新轮数 \(N\sim\mathrm{Geom}(p(\theta))\) 满足 \(\EE[N]=1/p(\theta)\)，随 \(\theta\) 严格递减，最小于 \(\theta=c\)。
为比较 \(\EE_x[T]\)，记在一次刷新时刻之后（即已产生候选 \(Y\sim\pi_\delta\)）的剩余步数为 \(\widetilde T\)。
在策略 \(a(y)=\theta/r_\delta(y)\) 下有递推
\[
\EE[\widetilde T]
=\sum_{y\in X}\pi_\delta(y)\Bigl(1-\frac{\theta}{r_\delta(y)}\Bigr)\Bigl(\EE[H_y]+\EE[\widetilde T]\Bigr),
\]
从而
\[
\EE[\widetilde T]
=\frac{\sum_y \pi_\delta(y)\bigl(1-\theta/r_\delta(y)\bigr)\,\EE[H_y]}{\sum_y \pi_\delta(y)\theta/r_\delta(y)}
=\frac{1}{\theta}-\frac{\sum_y \pi_\delta(y)/r_\delta(y)^2}{\sum_y \pi_\delta(y)/r_\delta(y)}.
\]
又由于首次刷新等待 \(H_x\) 与刷新候选序列独立且 \(\EE[H_x]=1/r_\delta(x)\)，有
\[
\EE_x[T]=\EE[H_x]+\EE[\widetilde T]
=\frac{1}{r_\delta(x)}+\frac{1}{\theta}-\frac{\sum_y \pi_\delta(y)/r_\delta(y)^2}{\sum_y \pi_\delta(y)/r_\delta(y)}.
\]
右端对 \(\theta\) 严格递减，故 \(\EE_x[T]\) 在 \((0,c]\) 上唯一最小于 \(\theta=c\)。
结合 \(\EE[N]\) 的严格单调性即得结论。证毕。
\end{proof}

\begin{theorem}[强平稳时刻的 \texorpdfstring{$\mathcal Q_\delta$}{Q}-有理母函数与分离距离上界]\label{thm:pom-diagonal-rate-sst-pgf-Qdelta}
沿用定理 \ref{thm:pom-diagonal-rate-refresh-accept-strong-stationary-time} 的最优策略 \(a(y)=c/r_\delta(y)\) 与强平稳时刻 \(T\)。
则对 \(|s|\) 充分小有严格恒等式
\[
\boxed{\;
\EE_x[s^{T}]
=\frac{s(t_x-\kappa_\delta)}{t_x-\kappa_\delta s}\cdot
\frac{c}{A_\delta^{2}}\cdot
\frac{1-s}{1-(1-c)s}\cdot
\frac{\kappa_\delta\,\mathcal P_\delta(\kappa_\delta s)}{\mathcal Q_\delta(\kappa_\delta s)}\; } ,
\]
其中 \(A_\delta=\sum_{x\in X}u_\delta(x)\) 且 \(A_\delta^{-2}=\sum_{y\in X} t_y/(t_y-\kappa_\delta)^2\)。
\par
并且对任意整数 \(m\ge 0\) 有强平稳时刻的一般上界
\[
\mathrm{sep}\!\bigl((K^\star_\delta)^m(x,\cdot),w\bigr)\le \PP_x(T>m),
\qquad
\bigl\|(K^\star_\delta)^m(x,\cdot)-w\bigr\|_{\mathrm{TV}}\le \PP_x(T>m),
\]
其中右端尾概率可由上述有理母函数经部分分式展开为有限几何和，从而给出可计算的混合剖面证书。
\end{theorem}
\begin{proof}
令 \(H_x\) 为起点 \(x\) 的首次刷新等待时间，则 \(\EE[s^{H_x}]=s(t_x-\kappa_\delta)/(t_x-\kappa_\delta s)\)（引理 \ref{lem:pom-diagonal-rate-refresh-holding-time-pgf} 与命题 \ref{prop:pom-diagonal-rate-diagonal-closure-rankone}）。
将 \(T\) 分解为首次刷新等待与“自刷新时刻起的剩余时间”之和，并对剩余时间母函数作一次刷新轮的递推（按接受/拒绝分解）得到
\[
F(s)=\frac{\sum_y \pi_\delta(y)a(y)}{1-\sum_y\pi_\delta(y)\bigl(1-a(y)\bigr)\EE[s^{H_y}]},
\qquad \EE_x[s^T]=\EE[s^{H_x}]\,F(s).
\]
在 \(a(y)=c/r_\delta(y)\) 下，代入 \(\pi_\delta(y)=1/(t_y-\kappa_\delta)\)、
\(
\EE[s^{H_y}]=s(t_y-\kappa_\delta)/(t_y-\kappa_\delta s)
\)
与恒等式
\[
1-s\sum_{y\in X}\frac{1}{t_y-\kappa_\delta s}=\frac{\mathcal Q_\delta(\kappa_\delta s)}{\kappa_\delta\mathcal P_\delta(\kappa_\delta s)},
\qquad
\sum_{y\in X}\frac{\pi_\delta(y)}{r_\delta(y)}=A_\delta^{-2},
\]
整理即得母函数闭式。
\par
分离距离与全变差的上界是强平稳时刻的通用性质；尾概率的有限几何和展开由有理母函数的部分分式分解给出。证毕。
\end{proof}

\paragraph{再生展开与刷新统计：更新过程、母函数有理性与 Laguerre--世俗 SCGF}
秩一刷新表示
\(
K^\star_\delta=\mathrm{diag}(1-r_\delta)+r_\delta\,\pi_\delta^{\mathsf T}
\)
不仅将谱与 resolvent 压缩为一维世俗方程（定理 \ref{thm:pom-maxent-markov-laguerre-secular-spectrum} 与命题 \ref{prop:pom-diagonal-rate-resolvent-entry-closed-form}），
亦把路径层面的随机性完全外置为一个更新过程：刷新候选为 i.i.d.，而停留长度为状态依赖几何分布。
这一再生表示在离散时间上给出刷新计数、再生奖励与熵率的闭式极限律，并在指数尺度上导出刷新计数的 SCGF，其 Perron 根由 Laguerre 变形的世俗多项式确定。

\begin{theorem}[秩一刷新核的再生展开表示]\label{thm:pom-diagonal-rate-refresh-regeneration-iid-expansion}
\leanverified{paper\_pom\_diagonal\_rate\_refresh\_regeneration\_iid\_expansion}
沿用命题 \ref{prop:pom-diagonal-rate-diagonal-closure-rankone} 的记号：
\(\delta\in(0,\delta_0)\)、\(\kappa:=\kappa_\delta>0\)、\(\{t_x\}_{x\in X}\subset(\kappa,\infty)\)，并令
\[
\pi_\delta(x):=\frac{1}{t_x-\kappa}\in\Delta(X),
\qquad
r_\delta(x):=\frac{t_x-\kappa}{t_x}\in(0,1).
\]
令 \(\{(Y_k,H_k)\}_{k\ge 1}\) 为 i.i.d.\ 随机对，满足对所有 \(y\in X\)、\(h\in\{1,2,\dots\}\) 有
\[
\PP(Y_k=y,\ H_k=h)=\pi_\delta(y)\,r_\delta(y)\,(1-r_\delta(y))^{h-1}.
\]
定义更新时刻 \(\sigma_0:=0\)、\(\sigma_k:=\sum_{j=1}^{k}H_j\)。
并定义离散时间过程 \((X_m)_{m\ge 0}\) 为
\[
X_m:=Y_k\qquad(\sigma_{k-1}\le m<\sigma_k).
\]
则：
\begin{enumerate}
  \item \((X_m)_{m\ge 0}\) 为以核 \(K^\star_\delta\) 为转移的马尔可夫链；
  \item \(\{\sigma_k\}_{k\ge 0}\) 为再生时刻：对每个 \(k\ge 0\)，在事件 \(\{m=\sigma_k\}\) 上，
  后继刷新候选 \(Y_{k+1}\) 与过去 \(\sigma\)-代数 \(\sigma(X_0,\dots,X_m)\) 独立且服从 \(\pi_\delta\)。
\end{enumerate}
\end{theorem}
\begin{proof}
由构造，若 \(\sigma_{k-1}\le m<\sigma_k-1\)，则 \(X_{m+1}=X_m=Y_k\)；
而在 \(m=\sigma_k-1\) 时，进入下一块并令 \(X_{m+1}=Y_{k+1}\)，其中 \(Y_{k+1}\sim\pi_\delta\) 且独立于过去。
因此对任意 \(x,y\in X\)，在条件 \(X_m=x\) 下，
\[
\PP(X_{m+1}=y\mid X_m=x)
=\PP(H_k\ge 2\mid Y_k=x)\mathbf 1_{\{y=x\}}
\;+\;\PP(H_k=1\mid Y_k=x)\pi_\delta(y)
=(1-r_\delta(x))\mathbf 1_{\{y=x\}}+r_\delta(x)\pi_\delta(y),
\]
这正是 \(K^\star_\delta\) 的秩一刷新表示（命题 \ref{prop:pom-diagonal-rate-diagonal-closure-rankone}），从而得到 (1)。
(2) 由 \(\{(Y_k,H_k)\}\) 的独立同分布性与块展开定义直接推出。证毕。
\end{proof}

\begin{proposition}[刷新间隔的有理母函数与 Hausdorff 完全单调尾]\label{prop:pom-diagonal-rate-refresh-interval-pgf-logderivative}
\leanverified{paper\_pom\_diagonal\_rate\_refresh\_interval\_pgf\_logderivative}
在定理 \ref{thm:pom-diagonal-rate-refresh-regeneration-iid-expansion} 的设定下，令 \(H:=H_1\) 为单块长度。
记全多项式
\[
\mathcal P_\delta(z):=\prod_{x\in X}(t_x-z).
\]
则对 \(|s|\) 充分小有严格恒等式
\begin{equation}\label{eq:pom-diagonal-rate-refresh-interval-pgf}
\boxed{\;
\Phi_H(s):=\EE[s^{H}]
=\sum_{x\in X}\pi_\delta(x)\frac{r_\delta(x)s}{1-(1-r_\delta(x))s}
=s\sum_{x\in X}\frac{1}{t_x-\kappa s}
\;=\;-\,s\,\frac{\mathcal P_\delta'(\kappa s)}{\mathcal P_\delta(\kappa s)}\; } .
\end{equation}
并且对任意整数 \(m\ge 0\)，尾概率满足
\begin{equation}\label{eq:pom-diagonal-rate-refresh-interval-tail}
\boxed{\;
\PP(H>m)=\sum_{x\in X}\pi_\delta(x)\,(1-r_\delta(x))^{m}
=\sum_{x\in X}\frac{1}{t_x-\kappa}\left(\frac{\kappa}{t_x}\right)^{m}\; } .
\end{equation}
因此 \(\{\PP(H>m)\}_{m\ge 0}\) 为 \([0,1)\) 上有限原子测度的 Hausdorff 矩序列，并满足完全单调性：
对任意 \(k,m\ge 0\)，有 \(({-1})^{k}(\Delta^{k}\PP(H>\cdot))_{m}\ge 0\)。
\end{proposition}
\begin{proof}
由条件分布 \(H\mid(Y=x)\sim\mathrm{Geom}(r_\delta(x))\)（支撑于 \(\{1,2,\dots\}\)）可得
\(
\EE[s^H]=\sum_x\pi_\delta(x)\frac{r_\delta(x)s}{1-(1-r_\delta(x))s}
\)。
代入 \(\pi_\delta(x)=1/(t_x-\kappa)\) 与 \(r_\delta(x)=(t_x-\kappa)/t_x\) 得
\(
\pi_\delta(x)\frac{r_\delta(x)s}{1-(1-r_\delta(x))s}
=s/(t_x-\kappa s)
\)，从而得到第二个等号。
最后用对数导数恒等式
\(
\sum_x(t_x-z)^{-1}=-\mathcal P_\delta'(z)/\mathcal P_\delta(z)
\)
得到 \eqref{eq:pom-diagonal-rate-refresh-interval-pgf} 的闭式。
\par
尾概率表示由 \(\PP(H>m\mid Y=x)=(1-r_\delta(x))^m\) 得到，并用 \(1-r_\delta(x)=\kappa/t_x\) 化简为 \eqref{eq:pom-diagonal-rate-refresh-interval-tail}。
完全单调性对 \(\lambda^m\) 逐项差分并用 \(\lambda\in(0,1)\) 即得。证毕。
\end{proof}

\begin{theorem}[刷新计数的更新极限定理]\label{thm:pom-diagonal-rate-refresh-count-renewal-lln-clt}
\leanverified{paper\_pom\_diagonal\_rate\_refresh\_count\_renewal\_lln\_clt}
在定理 \ref{thm:pom-diagonal-rate-refresh-regeneration-iid-expansion} 的设定下，令刷新计数
\[
N_m:=\max\{k:\ \sigma_k\le m\}\qquad(m\ge 0).
\]
则极限
\[
\EE[H]=\sum_{x\in X}\frac{\pi_\delta(x)}{r_\delta(x)}=\sum_{x\in X}\frac{t_x}{(t_x-\kappa)^2}
\]
有限，且几乎必然成立
\[
\frac{N_m}{m}\longrightarrow \frac{1}{\EE[H]}.
\]
若再记
\[
\EE[H^2]=\sum_{x\in X}\pi_\delta(x)\frac{2-r_\delta(x)}{r_\delta(x)^2}
=\sum_{x\in X}\frac{t_x(t_x+\kappa)}{(t_x-\kappa)^3},
\qquad
\Var(H)=\EE[H^2]-\EE[H]^2,
\]
则中心极限定律成立：
\[
\frac{N_m-\frac{m}{\EE[H]}}{\sqrt m}\ \Longrightarrow\ \mathcal N\!\left(0,\ \frac{\Var(H)}{\EE[H]^3}\right).
\]
\end{theorem}
\begin{proof}
由定理 \ref{thm:pom-diagonal-rate-refresh-regeneration-iid-expansion}，\((H_k)_{k\ge 1}\) 为 i.i.d.\ 且具有指数尾，故 \(\EE[H],\EE[H^2]<\infty\)。
更新过程的强大数定律与中心极限定理给出 \(N_m/m\to 1/\EE[H]\) 及所示归一化极限分布（标准更新计数定理）。证毕。
\end{proof}

\begin{theorem}[再生奖励的中心极限定理与显式方差]\label{thm:pom-diagonal-rate-refresh-reward-clt-variance}
\leanverified{paper\_pom\_diagonal\_rate\_refresh\_reward\_clt\_variance}
在定理 \ref{thm:pom-diagonal-rate-refresh-regeneration-iid-expansion} 的设定下，取任意可观测 \(f:X\to\RR\)，并令
\[
S_m(f):=\sum_{j=0}^{m-1}f(X_j).
\]
在平稳初始 \(X_0\sim w\) 下，令 \(\mu_f:=\sum_{x\in X}w(x)f(x)\)。
则存在方差常数 \(\sigma_f^2\ge 0\) 使
\[
\frac{S_m(f)-m\mu_f}{\sqrt m}\ \Longrightarrow\ \mathcal N(0,\sigma_f^2),
\qquad
\sigma_f^2=\frac{\Var\!\bigl(f(Y)H-\mu_f H\bigr)}{\EE[H]},
\]
其中 \((Y,H)\) 具有定理 \ref{thm:pom-diagonal-rate-refresh-regeneration-iid-expansion} 中的单块联合分布。
\end{theorem}
\begin{proof}
由块展开，令 \(R_k:=f(Y_k)H_k\) 为第 \(k\) 块奖励，则 \(R_k\) 为 i.i.d.\ 且 \(\EE[R_1^2]<\infty\)。
由再生奖励过程的标准中心极限定理，有
\(
\frac{S_m(f)-m\mu_f}{\sqrt m}\Rightarrow\mathcal N(0,\Var(R_1-\mu_f H_1)/\EE[H_1])
\)；
将 \((R_1,H_1)\) 写回为 \((f(Y)H,H)\) 即得所示方差形式。证毕。
\end{proof}

\begin{theorem}[刷新计数的 Laguerre--世俗 SCGF]\label{thm:pom-diagonal-rate-refresh-count-scgf-laguerre}
\leanverified{paper\_pom\_diagonal\_rate\_refresh\_count\_scgf\_laguerre}
在命题 \ref{prop:pom-diagonal-rate-diagonal-closure-rankone} 的秩一刷新表示下，对任意 \(\theta\in\RR\) 定义正核
\[
K_{\delta,\theta}(y\mid x):=(1-r_\delta(x))\mathbf 1_{\{y=x\}}+e^{\theta}r_\delta(x)\pi_\delta(y).
\]
则 \(K_{\delta,\theta}\) 的 Perron 根 \(\rho(\theta)>0\) 为方程
\begin{equation}\label{eq:pom-diagonal-rate-refresh-count-perron-equation}
\boxed{\;
\sum_{x\in X}\frac{1}{\rho(\theta)\,t_x-\kappa}=e^{-\theta}\;}
\end{equation}
的唯一正解。
等价地，令 \(z(\theta):=\kappa/\rho(\theta)\)，则 \(z(\theta)\in(0,\min_x t_x)\) 为 Laguerre 变形多项式
\begin{equation}\label{eq:pom-diagonal-rate-refresh-count-laguerre-deform}
\boxed{\;
\mathcal Q_{\delta,\theta}(z):=z\,\mathcal P_\delta'(z)+\kappa e^{-\theta}\,\mathcal P_\delta(z)\;}
\end{equation}
在该区间内的唯一零点，并满足 \(\rho(\theta)=\kappa/z(\theta)\)。
\par
在平稳初始下，刷新计数满足 SCGF 极限
\[
\psi(\theta):=\lim_{m\to\infty}\frac{1}{m}\log\EE_w\bigl[e^{\theta N_m}\bigr]=\log\rho(\theta),
\]
并且 \(\psi\) 在 \(\RR\) 上实解析且严格凸。
此外，
\[
\psi'(0)=\frac{1}{\EE[H]},
\qquad
\psi''(0)=\frac{\Var(H)}{\EE[H]^3},
\]
其中 \(\EE[H],\Var(H)\) 由定理 \ref{thm:pom-diagonal-rate-refresh-count-renewal-lln-clt} 的闭式给出。
\end{theorem}
\begin{proof}
由秩一表示，\(K_{\delta,\theta}=\mathrm{diag}(1-r_\delta)+e^{\theta}r_\delta\pi_\delta^{\mathsf T}\)。
对 \(\rho I-K_{\delta,\theta}\) 应用秩一行列式引理并用
\(
\pi_\delta(x)r_\delta(x)=1/t_x
\)
得到其世俗方程等价于 \eqref{eq:pom-diagonal-rate-refresh-count-perron-equation}，从而 Perron 根的唯一性由左端对 \(\rho\) 的严格单调性与正性给出。
令 \(z=\kappa/\rho\) 并化简可得 \(\mathcal Q_{\delta,\theta}(z)=0\)，即 \eqref{eq:pom-diagonal-rate-refresh-count-laguerre-deform}。
\par
对 SCGF，记
\(
F_m(\theta):=\EE_w[e^{\theta N_m}]
\)
并将刷新事件视为沿路径的乘法权重，则由 Feynman--Kac 迭代有
\(
F_m(\theta)=w^{\mathsf T}K_{\delta,\theta}^{\,m}\mathbf 1
\)。
由于 \(K_{\delta,\theta}\) 为逐项正矩阵，其幂增长由 Perron 根控制，故
\(
\lim_{m\to\infty}\frac{1}{m}\log F_m(\theta)=\log\rho(\theta)
\)。
解析性与严格凸性由 Perron 根对参数的解析扰动与 \(\log\rho(\theta)\) 的严格对数凸性给出。
最后，\(\psi'(0)\) 与 \(\psi''(0)\) 分别等于 \(\lim_m \EE[N_m]/m\) 与 \(\lim_m\Var(N_m)/m\)，而这两者由更新计数定理给出为 \(1/\EE[H]\) 与 \(\Var(H)/\EE[H]^3\)（定理 \ref{thm:pom-diagonal-rate-refresh-count-renewal-lln-clt}）。证毕。
\end{proof}

\begin{proposition}[熵率的再生编码闭式]\label{prop:pom-diagonal-rate-refresh-entropy-rate-renewal-coding}
沿用定理 \ref{thm:pom-diagonal-rate-refresh-regeneration-iid-expansion} 的记号，令
\(
\mathsf h:=h(K^\star_\delta)
\)
为定理 \ref{thm:pom-maxent-markov-unique-optimal-kernel} 所定义之熵率。
则有再生编码恒等式
\[
\boxed{\;
\mathsf h=\frac{H(Y,H)}{\EE[H]}\;}
\]
并且可写成完全闭式
\[
\mathsf h
=\frac{H(\pi_\delta)+\EE_{\pi_\delta}\!\left[-\log r_\delta(Y)-\frac{1-r_\delta(Y)}{r_\delta(Y)}\log(1-r_\delta(Y))\right]}{\sum_{x\in X}\pi_\delta(x)/r_\delta(x)}.
\]
\end{proposition}
\begin{proof}
由再生展开，路径由独立块 \((Y_k,H_k)\) 编码，且每块在时间尺度上占据 \(H_k\) 个符号。
对平稳过程应用 Shannon--McMillan--Breiman 定理（有限状态平稳遍历链），单位时间的信息率等于单位块信息量与单位块长度之比：
\(
\mathsf h=H(Y,H)/\EE[H]
\)。
又 \(H(Y,H)=H(\pi_\delta)+\EE_{\pi_\delta}[H(H\mid Y)]\)，并且对 \(\mathrm{Geom}(r)\)（支撑 \(\{1,2,\dots\}\)）有熵闭式
\(
H(\mathrm{Geom}(r))=-\log r-\frac{1-r}{r}\log(1-r)
\)，
代回即得所示表达。证毕。
\end{proof}

\endinput


\paragraph{再生编码、更新极限定理与刷新计数的 Laguerre 世俗 SCGF}
固定 \(\delta\in(0,\delta_0)\) 并沿用命题 \ref{prop:pom-diagonal-rate-diagonal-closure-rankone} 与定理 \ref{thm:pom-diagonal-rate-sturm-cauchy-eigenbasis} 的记号。
为简洁起见以下记 \(\kappa:=\kappa_\delta\)、\(t_x\) 不变，并令
\[
r(x):=\frac{t_x-\kappa}{t_x}\in(0,1),
\qquad
\pi(x):=\frac{1}{t_x-\kappa}\in\Delta(X),
\qquad
K:=K^\star_\delta.
\]
则 \(K\) 具有停留--刷新表示
\begin{equation}\label{eq:pom-diagonal-rate-refresh-stay-refresh}
K(y\mid x)=(1-r(x))\,\mathbf 1_{\{y=x\}}+r(x)\pi(y),
\qquad x,y\in X.
\end{equation}

\begin{theorem}[再生编码：\texorpdfstring{$(Y_k,H_k)$}{(Y,H)} 的 i.i.d.\ 展开生成停留--刷新链]\label{thm:pom-diagonal-rate-refresh-renewal-coding}
令 \((Y_k,H_k)_{k\ge 1}\) 为 i.i.d.\ 随机对，满足对任意 \(y\in X\)、\(h\in\{1,2,\dots\}\),
\[
\PP(Y_1=y,\ H_1=h)=\pi(y)\,r(y)\,(1-r(y))^{h-1}.
\]
记再生时刻 \(\sigma_0:=0\)、\(\sigma_k:=\sum_{j=1}^{k}H_j\)。
定义展开过程 \((X_m)_{m\ge 0}\) 为
\[
X_m:=Y_k\qquad(\sigma_{k-1}\le m<\sigma_k).
\]
则 \((X_m)\) 为以核 \(K\) 转移、初始分布为 \(\pi\) 的马尔可夫链。
并且每个 \(\sigma_k\) 都是刷新时刻：\((X_{\sigma_k})_{k\ge 0}\) 为 i.i.d.\ 且边缘分布为 \(\pi\)，并与 \(\{H_k\}\) 的未来独立。
\end{theorem}
\begin{proof}
由构造，\(X_m\) 在每个块内常值；因此只需计算“从 \(x\) 出发的一步转移”。
当 \(X_m=x\) 时，事件 \(\{X_{m+1}=x\}\) 等价于“当前块尚未结束”。
由几何分布的无记忆性，
\[
\PP(X_{m+1}=x\mid X_m=x)=\PP(H_1\ge 2\mid H_1\ge 1,\ Y_1=x)=1-r(x).
\]
而事件 \(\{X_{m+1}=y,\ y\ne x\}\) 等价于“当前步结束块并刷新到 \(y\)”；
同样由无记忆性与刷新采样 \(Y_{k+1}\sim\pi\)，有
\[
\PP(X_{m+1}=y\mid X_m=x)=r(x)\pi(y)\qquad(y\in X).
\]
综上一步转移即为 \eqref{eq:pom-diagonal-rate-refresh-stay-refresh}，从而 \((X_m)\) 为 \(K\)-链。
刷新时刻的 i.i.d.\ 性与独立性由 \((Y_k,H_k)\) 的 i.i.d.\ 与展开映射的确定性直接给出。证毕。
\end{proof}
\leanverified{paper\_pom\_diagonal\_rate\_refresh\_renewal\_coding}

\begin{corollary}[平稳分布的长度偏置形式与 \texorpdfstring{$t$}{t}-闭式]\label{cor:pom-diagonal-rate-refresh-length-biased-stationary}
记
\[
\EE[H]:=\sum_{x\in X}\frac{\pi(x)}{r(x)}=\sum_{x\in X}\frac{t_x}{(t_x-\kappa)^2}\in(1,\infty).
\]
则 \(K\) 的唯一平稳分布满足长度偏置恒等式
\[
\boxed{\;
w(x)=\frac{\pi(x)\,\EE[H\mid Y=x]}{\EE[H]}
=\frac{\pi(x)/r(x)}{\sum_{u\in X}\pi(u)/r(u)}
=\frac{\dfrac{t_x}{(t_x-\kappa)^2}}{\sum_{u\in X}\dfrac{t_u}{(t_u-\kappa)^2}}\; } .
\]
\end{corollary}
\begin{proof}
由 \(r(x)\pi(y)>0\)（\(\pi\) 全支撑且 \(r(x)\in(0,1)\)），链不可约且非周期，故平稳分布唯一。
由定理 \ref{thm:pom-diagonal-rate-refresh-renewal-coding} 的再生编码，每个块的落点为 \(Y\sim\pi\)，并在该落点停留 \(H\mid(Y=x)\sim\mathrm{Geom}(r(x))\)，从而 \(\EE[H\mid Y=x]=1/r(x)\)。
再生链的占时比例等于“单位块内期望占时”与“单位块期望长度”之比，故
\(
w(x)\propto \pi(x)\EE[H\mid Y=x]=\pi(x)/r(x)
\)；
归一化即得结论。证毕。
\end{proof}
\leanverified{paper\_pom\_diagonal\_rate\_refresh\_length\_biased\_stationary}

\begin{proposition}[刷新间隔的 \texorpdfstring{$\mathcal P_\delta$}{P}-有理母函数与 Hausdorff 完全单调尾]\label{prop:pom-diagonal-rate-refresh-holding-interval-pgf-hausdorff}
令 \(H\) 表示以 \(Y\sim\pi\) 为块落点时的单块长度（即 \(H\overset{d}{=}H_1\)）。
则对 \(|s|\) 充分小有严格恒等式
\[
\boxed{\;
\Phi_H(s):=\EE[s^H]
=\sum_{x\in X}\pi(x)\frac{r(x)s}{1-(1-r(x))s}
=s\sum_{x\in X}\frac{1}{t_x-\kappa s}
=-\,s\,\frac{\mathcal P_\delta'(\kappa s)}{\mathcal P_\delta(\kappa s)}\; } .
\]
并且对任意整数 \(m\ge 0\) 有尾概率闭式
\[
\boxed{\;
\PP(H>m)=\sum_{x\in X}\pi(x)\,(1-r(x))^{m}
=\sum_{x\in X}\frac{1}{t_x-\kappa}\left(\frac{\kappa}{t_x}\right)^{m}\; } .
\]
因此 \(\{\PP(H>m)\}_{m\ge 0}\) 为 \([0,1)\) 上有限原子测度的 Hausdorff 矩序列，并满足完全单调性
\[
(-1)^k(\Delta^k\PP(H>\cdot))_m\ge 0\qquad(k,m\ge 0).
\]
\end{proposition}
\leanverified{paper\_pom\_diagonal\_rate\_refresh\_holding\_interval\_pgf\_hausdorff}
\begin{proof}
由 \(\PP(H=h\mid Y=x)=r(x)(1-r(x))^{h-1}\)（\(h\ge 1\)）直接求和得
\(
\EE[s^H\mid Y=x]=\frac{r(x)s}{1-(1-r(x))s}
\)，
再对 \(Y\sim\pi\) 取期望得到第一条母函数表示。
由 \(r(x)\pi(x)=1/t_x\) 与 \(1-(1-r(x))s=(t_x-\kappa s)/t_x\) 化简得
\(
\Phi_H(s)=s\sum_x (t_x-\kappa s)^{-1}
\)。
最后，由 \(\mathcal P_\delta(z)=\prod_x(t_x-z)\) 的对数导数恒等式
\(
\sum_x (t_x-z)^{-1}=-\mathcal P_\delta'(z)/\mathcal P_\delta(z)
\)
得到有理闭式。
\par
尾概率由 \(\PP(H>m\mid Y=x)=(1-r(x))^m\) 与 \(1-r(x)=\kappa/t_x\) 得到。
Hausdorff 与完全单调性由离散测度
\(
\nu:=\sum_x \pi(x)\delta_{\kappa/t_x}
\)
满足
\(
\PP(H>m)=\int_0^1 \lambda^m\,\nu(d\lambda)
\)
直接推出。证毕。
\end{proof}

\begin{theorem}[刷新计数的更新 LLN 与 CLT]\label{thm:pom-diagonal-rate-refresh-count-renewal-clt}
\leanverified{paper\_pom\_diagonal\_rate\_refresh\_count\_renewal\_clt}
沿用定理 \ref{thm:pom-diagonal-rate-refresh-renewal-coding} 的再生时刻 \((\sigma_k)\)，并定义刷新计数
\[
N_m:=\max\{k:\ \sigma_k\le m\}\qquad(m\ge 0).
\]
则几乎必然有
\[
\frac{N_m}{m}\longrightarrow \frac{1}{\EE[H]}.
\]
并且当 \(\EE[H^2]<\infty\)（在本核上必然成立）时有中心极限定理
\[
\frac{N_m-\frac{m}{\EE[H]}}{\sqrt m}\ \Longrightarrow\ \mathcal N\!\left(0,\ \frac{\Var(H)}{\EE[H]^3}\right).
\]
其中
\[
\EE[H^2]=\sum_{x\in X}\pi(x)\,\frac{2-r(x)}{r(x)^2}
=\sum_{x\in X}\frac{t_x(t_x+\kappa)}{(t_x-\kappa)^3}.
\]
\end{theorem}
\begin{proof}[证明要点]
由定理 \ref{thm:pom-diagonal-rate-refresh-renewal-coding}，\((H_k)_{k\ge 1}\) 为 i.i.d.\ 且与刷新计数满足标准更新关系 \(N_m=\max\{k:\sum_{j\le k}H_j\le m\}\)。
因此大数定律与中心极限定理为经典更新过程定理的直接推论。
二阶矩闭式由 \(\mathrm{Geom}(p)\)（支撑 \(\{1,2,\dots\}\)）满足 \(\EE[H^2]=(2-p)/p^2\) 并对 \(p=r(Y)\) 条件化求期望得到；再代入 \(r(x)=(t_x-\kappa)/t_x\)、\(\pi(x)=1/(t_x-\kappa)\) 化简即得。证毕。
\end{proof}

\begin{theorem}[再生奖励 CLT：一般加性泛函的显式方差]\label{thm:pom-diagonal-rate-refresh-reward-clt}
\leanverified{paper\_pom\_diagonal\_rate\_refresh\_reward\_clt}
对任意可观测 \(f:X\to\RR\)，记
\[
S_m(f):=\sum_{j=0}^{m-1}f(X_j),
\qquad
\mu_f:=\sum_{x\in X}w(x)f(x).
\]
则在平稳初始 \(X_0\sim w\) 下有
\[
\frac{S_m(f)-m\mu_f}{\sqrt m}\ \Longrightarrow\ \mathcal N(0,\sigma_f^2),
\qquad
\sigma_f^2=\frac{\Var\!\bigl(H(f(Y)-\mu_f)\bigr)}{\EE[H]},
\]
其中 \((Y,H)\) 的联合分布如定理 \ref{thm:pom-diagonal-rate-refresh-renewal-coding}。
\end{theorem}
\begin{proof}[证明要点]
由再生编码，单块奖励为 \(\sum_{j=\sigma_{k-1}}^{\sigma_k-1} f(X_j)=H_k f(Y_k)\)；
并且在平稳初始下中心化后单块奖励可写为
\(
H_k(f(Y_k)-\mu_f)
\)，其均值为 \(0\) 且二阶矩有限。
由更新奖励过程的中心极限定理得到所示结论与方差表达。证毕。
\end{proof}

\begin{theorem}[刷新计数的倾斜算子与 Laguerre 世俗 Perron 根]\label{thm:pom-diagonal-rate-refresh-count-tilt-perron-secular}
\leanverified{paper\_pom\_diagonal\_rate\_refresh\_count\_tilt\_perron\_secular}
对 \(\theta\in\RR\) 定义正算子
\[
K_\theta(y\mid x):=(1-r(x))\,\mathbf 1_{\{y=x\}}+e^{\theta}r(x)\pi(y).
\]
则 \(K_\theta\) 的 Perron 根 \(\rho(\theta)>0\) 是方程
\[
\boxed{\;
\sum_{x\in X}\frac{1}{\rho(\theta)\,t_x-\kappa}=e^{-\theta}\;}
\]
的唯一正解。
等价地，令 \(z(\theta):=\kappa/\rho(\theta)\)，则 \(z(\theta)\in(0,\min_x t_x)\) 为多项式
\[
\mathcal Q_{\delta,\theta}(z):=z\,\mathcal P_\delta'(z)+\kappa e^{-\theta}\,\mathcal P_\delta(z)
\]
在 \((0,\min_x t_x)\) 内的唯一零点，并且 \(\rho(\theta)=\kappa/z(\theta)\)。
\end{theorem}
\begin{proof}
由 \eqref{eq:pom-diagonal-rate-refresh-stay-refresh}，\(K_\theta\) 仍为“对角 \(+\) 秩一”算子：
\[
K_\theta=\mathrm{diag}\!\Bigl(\frac{\kappa}{t_x}\Bigr)_{x\in X}+e^{\theta}\,r\,\pi^{\mathsf T}.
\]
对特征方程 \(\det(\rho I-K_\theta)=0\) 应用秩一行列式引理并用 \(\pi(x)r(x)=1/t_x\)，得到
\[
1=e^{\theta}\sum_{x\in X}\frac{\pi(x)r(x)}{\rho-\kappa/t_x}
=e^{\theta}\sum_{x\in X}\frac{1}{\rho t_x-\kappa},
\]
即所示标量方程。
函数 \(\rho\mapsto \sum_x(\rho t_x-\kappa)^{-1}\) 在 \((\kappa/\min t_x,\infty)\) 上连续严格递减且取值从 \(+\infty\) 到 \(0\)，故对任意 \(\theta\) 存在唯一正解 \(\rho(\theta)\)，并由 Perron--Frobenius 与正性可知该解即为谱半径。
\par
再令 \(z=\kappa/\rho\)，则标量方程等价于
\(
\sum_x (t_x-z)^{-1}=\kappa e^{-\theta}/z
\)。
将其改写为对数导数形式并乘以 \(\mathcal P_\delta(z)\) 即得 \(\mathcal Q_{\delta,\theta}(z)=0\)；
唯一性由同一单调性得到。证毕。
\end{proof}

\begin{corollary}[刷新计数的 SCGF 与二阶导数闭式]\label{cor:pom-diagonal-rate-refresh-count-scgf-derivatives}
\leanverified{paper\_pom\_diagonal\_rate\_refresh\_count\_scgf\_derivatives}
在平稳初始 \(X_0\sim w\) 下，刷新计数满足极限
\[
\psi(\theta):=\lim_{m\to\infty}\frac{1}{m}\log \EE_w\!\bigl[e^{\theta N_m}\bigr]=\log\rho(\theta),
\]
其中 \(\rho(\theta)\) 为定理 \ref{thm:pom-diagonal-rate-refresh-count-tilt-perron-secular} 的 Perron 根。
并且 \(\psi\) 在 \(\theta=0\) 的邻域内实解析，且导数满足
\[
\psi'(0)=\frac{1}{\EE[H]},
\qquad
\psi''(0)=\frac{\Var(H)}{\EE[H]^3}.
\]
写
\[
T_2:=\sum_{x\in X}\frac{t_x}{(t_x-\kappa)^2}=\EE[H],
\qquad
U_3:=\sum_{x\in X}\frac{t_x^{2}}{(t_x-\kappa)^3},
\]
则
\[
\boxed{\;
\psi''(0)=\frac{1}{T_2}\left(\frac{2U_3}{T_2^2}-1\right)-\frac{1}{T_2^2}\; } .
\]
\end{corollary}
\begin{proof}[证明要点]
由 Feynman--Kac 展开，\(\EE_x[e^{\theta N_m}]=\sum_y (K_\theta^m)(y\mid x)\)，故增长率由谱半径控制并给出 \(\psi(\theta)=\log\rho(\theta)\)。
对定理 \ref{thm:pom-diagonal-rate-refresh-count-tilt-perron-secular} 的标量方程在 \(\theta=0\) 处作隐函数微分，得到
\(
\rho'(0)=1/T_2
\)
与
\(
\rho''(0)=(2U_3/T_2^2-1)/T_2
\)；
再用 \(\psi=\log\rho\) 的恒等式 \(\psi''=\rho''/\rho-(\rho'/\rho)^2\) 并代入 \(\rho(0)=1\) 即得所示闭式。
最后一条 \(\psi''(0)=\Var(H)/\EE[H]^3\) 与定理 \ref{thm:pom-diagonal-rate-refresh-count-renewal-clt} 的方差常数一致。证毕。
\end{proof}

\begin{proposition}[熵率的再生编码公式]\label{prop:pom-diagonal-rate-refresh-entropy-rate-renewal-coding-secondary}
记马尔可夫链熵率
\[
h(K)=\sum_{x\in X} w(x)\,H\!\bigl(K(\cdot\mid x)\bigr).
\]
则沿用定理 \ref{thm:pom-diagonal-rate-refresh-renewal-coding} 的 \((Y,H)\)，有严格恒等式
\[
\boxed{\;
h(K)=\frac{H(Y,H)}{\EE[H]}\; } .
\]
\end{proposition}
\begin{proof}[证明要点]
定理 \ref{thm:pom-diagonal-rate-refresh-renewal-coding} 给出一个从 i.i.d.\ 块序列 \((Y_k,H_k)\) 到路径 \((X_m)\) 的确定性展开映射。
块间独立性给出 \(H(Y_1,\dots,Y_n,H_1,\dots,H_n)=n\,H(Y,H)\)。
另一方面，展开后前 \(m\) 步路径信息量由包含的块数 \(N_m\) 控制，并满足
\(
m^{-1}H(X_0,\dots,X_{m-1})\to H(Y,H)/\EE[H]
\)；
而平稳马尔可夫链的熵率等于该极限，故结论成立。证毕。
\end{proof}

\endinput



\endinput


\begin{theorem}[对角统计的完备性：由保持率向量回收 \texorpdfstring{$(\kappa_\delta,\delta,w,K^\star_\delta)$}{(kappa,delta,w,K*)}]\label{thm:pom-diagonal-rate-diagonal-statistics-complete}
在命题 \ref{prop:pom-diagonal-rate-diagonal-closure-rankone} 的设定下，设仅给定最优核对角
\(
p_\delta(x)=K^\star_\delta(x\mid x)
\)。
定义
\[
F(\kappa):=\sum_{x\in X}\frac{p_\delta(x)}{1+\kappa-\kappa p_\delta(x)}.
\]
则成立：
\begin{enumerate}
  \item \textbf{（\(\kappa_\delta\) 的唯一恢复）}
  方程 \(F(\kappa)=1\) 在 \((0,\infty)\) 上存在唯一解；该唯一解等于最优对偶参数 \(\kappa_\delta\)。

  \item \textbf{（背景分布的代数闭式）}
  令
  \[
  \pi_\delta(x):=\frac{p_\delta(x)}{1+\kappa_\delta-\kappa_\delta p_\delta(x)}\qquad(x\in X),
  \]
  则 \(\pi_\delta\in\Delta(X)\)，且与结论 \ref{cor:pom-maxent-markov-optimal-coupling-latent-binary-mixture} 的背景分布一致。
  记 \(s_2:=\sum_x \pi_\delta(x)^2\)，则失真由对角唯一确定为
  \[
  \delta=\frac{1-s_2}{1+\kappa_\delta s_2},
  \qquad\text{等价地}\qquad
  s_2=\frac{1-\delta}{1+\kappa_\delta\delta}.
  \]

  \item \textbf{（平稳分布的对角闭式）}
  平稳分布由对角显式回收：
  \[
  w(x)=\frac{\dfrac{p_\delta(x)}{(1+\kappa_\delta-\kappa_\delta p_\delta(x))^{2}}}{\sum_{y\in X}\dfrac{p_\delta(y)}{(1+\kappa_\delta-\kappa_\delta p_\delta(y))^{2}}}\qquad(x\in X).
  \]

  \item \textbf{（全核的对角闭式）}
  对任意 \(x,y\in X\)，
  \[
  K^\star_\delta(y\mid x)=
  \begin{cases}
  p_\delta(x),&y=x,\\[6pt]
  \Bigl(1-\dfrac{\kappa_\delta}{1+\kappa_\delta}p_\delta(x)\Bigr)\,\pi_\delta(y),&y\ne x.
  \end{cases}
  \]
\end{enumerate}
\end{theorem}
\begin{proof}[证明要点]
由命题 \ref{prop:pom-diagonal-rate-diagonal-closure-rankone}，\(p_\delta(x)=(1+\kappa_\delta)/t_x\) 且
\(
\sum_x (t_x-\kappa_\delta)^{-1}=1
\)。
注意
\(
t_x=(1+\kappa_\delta)/p_\delta(x)
\)，从而
\[
\frac{1}{t_x-\kappa_\delta}
=\frac{1}{\frac{1+\kappa_\delta}{p_\delta(x)}-\kappa_\delta}
=\frac{p_\delta(x)}{1+\kappa_\delta-\kappa_\delta p_\delta(x)}.
\]
将其对 \(x\) 求和即得 \(F(\kappa_\delta)=1\)。
另一方面，对任意固定的 \(p_\delta(x)\in(0,1)\)，函数 \(\kappa\mapsto p_\delta(x)/(1+\kappa-\kappa p_\delta(x))\) 在 \([0,\infty)\) 上连续严格递减，
从而 \(F\) 在 \([0,\infty)\) 上连续严格递减且 \(F(\infty)=0\)。
由结论 \ref{cor:pom-maxent-markov-no-oscillation-endpoints}，\(\mathrm{Spec}(K^\star_\delta)\subseteq[0,1]\) 且 \(\lambda_1=1\)，故 \(\sum_x p_\delta(x)=\mathrm{tr}(K^\star_\delta)\ge 1\)；
并且当 \(\delta\in(0,\delta_0)\) 时 \(K^\star_\delta\) 不可能为秩一核（否则必退化为独立核并迫使 \(\delta=\delta_0\)），因此 \(\sum_x p_\delta(x)>1\)，即 \(F(0)>1\)。
由介值定理与严格单调性得到 \((0,\infty)\) 上唯一根，且必等于 \(\kappa_\delta\)。
\par
\(\pi_\delta\in\Delta(X)\) 由 \(F(\kappa_\delta)=1\) 直接推出；并由上式知其等于 \(1/(t_x-\kappa_\delta)=u_\delta(x)/A_\delta\)，故与背景分布一致。
写 \(u_\delta(x)=A_\delta\pi_\delta(x)\) 并代入归一化恒等式
\(
1=A_\delta^2+\kappa_\delta\sum_x u_\delta(x)^2
=A_\delta^2\bigl(1+\kappa_\delta s_2\bigr)
\)
得 \(A_\delta^2=(1+\kappa_\delta s_2)^{-1}\)。
再由对角质量
\(
1-\delta=\sum_x P^\star_\delta(x,x)=(1+\kappa_\delta)\sum_x u_\delta(x)^2=(1+\kappa_\delta)A_\delta^2 s_2
\)
即可解出 \(\delta=(1-s_2)/(1+\kappa_\delta s_2)\)。
\par
由 \(w(x)=t_xu_\delta(x)^2\) 与 \(u_\delta(x)=A_\delta\pi_\delta(x)=A_\delta/(t_x-\kappa_\delta)\) 得
\[
w(x)=A_\delta^2\,\frac{t_x}{(t_x-\kappa_\delta)^2}.
\]
将 \(t_x=(1+\kappa_\delta)/p_\delta(x)\) 代入并用 \(\sum_x w(x)=1\) 归一化，即得到 \(w\) 的对角闭式。
全核闭式由命题 \ref{prop:pom-diagonal-rate-diagonal-closure-rankone} 将 \(t_x\) 与 \((t_x-\kappa_\delta)^{-1}=\pi_\delta(x)\) 代回即可。证毕。
\par
可复现核验脚本：\texttt{scripts/exp\_pom\_diagonal\_rate\_diagonal\_observation\_audit.py}。
\end{proof}
\leanverified{paper_pom_diagonal_rate_diagonal_statistics_complete}

\begin{corollary}[对角 Laguerre 变换：特征行列式的闭式]\label{cor:pom-diagonal-rate-diagonal-laguerre-determinant}
在定理 \ref{thm:pom-diagonal-rate-diagonal-statistics-complete} 的设定下，定义对角多项式
\[
\widetilde{\mathcal P}_\delta(z):=\prod_{x\in X}\Bigl(1-\frac{p_\delta(x)}{1+\kappa_\delta}z\Bigr).
\]
则对任意 \(z\in\CC\) 有严格恒等式
\[
\det(I-zK^\star_\delta)=\widetilde{\mathcal P}_\delta(\kappa_\delta z)+z\,\widetilde{\mathcal P}_\delta'(\kappa_\delta z).
\]
\end{corollary}
\begin{proof}[证明要点]
由 \(p_\delta(x)=(1+\kappa_\delta)/t_x\) 得
\(
\mathcal P_\delta(z)=\prod_x(t_x-z)=\mathcal P_\delta(0)\widetilde{\mathcal P}_\delta(z)
\)。
将其代入
\(
\mathcal Q_\delta(z)=z\mathcal P_\delta'(z)+\kappa_\delta\mathcal P_\delta(z)
\)
与推论 \ref{cor:pom-maxent-markov-zeta-determinant-single-polynomial} 的恒等式
\(
\det(I-zK^\star_\delta)=\mathcal Q_\delta(\kappa_\delta z)/(\kappa_\delta\mathcal P_\delta(0))
\)
并约去公共因子即可。证毕。
\end{proof}

\begin{corollary}[对角即全谱：特征值与 \texorpdfstring{$\zeta_\delta$}{zeta} 极点的单次恢复]\label{cor:pom-diagonal-rate-diagonal-determines-spectrum}
在推论 \ref{cor:pom-diagonal-rate-diagonal-laguerre-determinant} 的设定下，
\(\zeta_\delta(z)=\det(I-zK^\star_\delta)^{-1}\) 的全部极点（含重数）由方程
\[
\widetilde{\mathcal P}_\delta(\kappa_\delta z)+z\,\widetilde{\mathcal P}_\delta'(\kappa_\delta z)=0
\]
的全部零点给出；等价地，最优核的特征值多重集由上述零点的倒数（再加上平凡特征值 \(1\)）完全确定。
\end{corollary}

\begin{corollary}[速率—对角 Euler 恒等式]\label{cor:pom-diagonal-rate-euler-diagonal-identity}
取随机变量 \(X\sim w\)，并记 \(p_\delta(X):=K^\star_\delta(X\mid X)\)。
则对任意 \(\delta\in(0,\delta_0)\) 有严格恒等式
\[
R_w(\delta)=H(w)+\EE_w\!\bigl[\log p_\delta(X)\bigr]+\delta\,R_w'(\delta).
\]
等价地，若记 \(\kappa_\delta=\exp\!\bigl(-R_w'(\delta)\bigr)-1\)，则
\[
R_w(\delta)=H(w)+\EE_w\!\bigl[\log p_\delta(X)\bigr]-\delta\log(1+\kappa_\delta).
\]
\end{corollary}
\begin{proof}[证明要点]
由命题 \ref{prop:pom-diagonal-rate-diagonal-closure-rankone} 与 \(w(x)=t_xu_\delta(x)^2\) 得
\(
p_\delta(x)=(1+\kappa_\delta)u_\delta(x)^2/w(x)
\)，从而
\[
\EE_w[\log p_\delta(X)]
=\log(1+\kappa_\delta)+2\sum_x w(x)\log u_\delta(x)-\sum_x w(x)\log w(x).
\]
将其与推论 \ref{cor:pom-diagonal-rate-scalar-collapse-rate} 的闭式速率表达式合并即可得到
\(
R_w(\delta)=H(w)+\EE_w[\log p_\delta(X)]-\delta\log(1+\kappa_\delta)
\)。
再由对偶敏感性公式 \(R_w'(\delta)=-\log(1+\kappa_\delta)\)（定理 \ref{thm:pom-diagonal-rate-explicit-param-and-derivative}）得到等价写法。证毕。
\end{proof}

\begin{proposition}[对角对数的回馈律：\texorpdfstring{$\delta$}{delta}-导数由曲率给出]\label{prop:pom-diagonal-rate-diagonal-log-feedback}
在 \(\delta\in(0,\delta_0)\) 上，函数 \(\delta\mapsto \EE_w[\log p_\delta(X)]\) 可微且满足
\[
\frac{\mathrm d}{\mathrm d\delta}\,\EE_w[\log p_\delta(X)]
=-\delta\,R_w''(\delta).
\]
特别地，由命题 \ref{prop:pom-diagonal-rate-analytic-strict-convexity-curvature} 的严格凸性 \(R_w''(\delta)>0\) 可知
\(\EE_w[\log p_\delta(X)]\) 在 \((0,\delta_0)\) 上严格递减。
并且结合其曲率表示可等价写为
\[
\frac{\mathrm d}{\mathrm d\delta}\,\EE_w[\log p_\delta(X)]
=-\frac{\delta}{\Var_{P^\star_\delta}(Z)},
\qquad
Z:=\mathbf 1_{\{X=Y\}},\ (X,Y)\sim P^\star_\delta.
\]
\leanpartial{paper\_pom\_diagonal\_rate\_diagonal\_log\_feedback}{Lean signature omits the hypotheses needed to derive the derivative identity for \texttt{paper\_pom\_diagonal\_rate\_diagonal\_log\_feedback}.}
\end{proposition}
\begin{proof}
对推论 \ref{cor:pom-diagonal-rate-euler-diagonal-identity} 的恒等式两边对 \(\delta\) 求导，并消去公共项 \(R_w'(\delta)\) 即得
\(
\frac{\mathrm d}{\mathrm d\delta}\EE_w[\log p_\delta(X)]=-\delta R_w''(\delta)
\)。
最后一式由命题 \ref{prop:pom-diagonal-rate-analytic-strict-convexity-curvature} 的曲率表示 \(R_w''(\delta)=1/\Var_{P^\star_\delta}(Z)\) 代入得到。证毕。
\end{proof}
